% =============================================================================
% BAB I PENDAHULUAN -- teks panduan templat Skripsi Sain-Tek-Kes 2026
% Ganti seluruh teks panduan dengan isi skripsi.
% =============================================================================

\chapter{PENDAHULUAN}

Bab pendahuluan memuat latar belakang atau justifikasi dipilihnya topik karya
ilmiah tugas akhir, perumusan atau pendekatan penyelesaian masalah, tujuan,
manfaat, dan ruang lingkup. Di dalam pendahuluan dijelaskan perumusan atau
pendekatan penyelesaian masalah dan alasan pemilihan metode yang digunakan.
Merujuk pada proses perumusan masalah tugas akhir, bagian Kerangka Pikir dan
Hipotesis dapat ditulis di sini, tidak ditulis dalam bab tersendiri. Untuk
membantu mengikuti alur pikir secara skematis, dapat juga dibuat bagan alir
kerangka proses dan rumusan masalah serta pencapaian tujuan tugas akhir.

\section{Latar Belakang}

Latar belakang berisi penjelasan alasan memilih topik dan pentingnya kajian
tugas akhir itu dilakukan berdasarkan alasan teoretis dan praktis, serta
bagaimana masalah tersebut dapat diatasi dan manfaat dari penyelesaian masalah.
Paparan tidak berbelit-belit atau dimulai dengan latar belakang yang umum.
Pernyataan mengenai apa yang diteliti dan apa yang diharapkan diawali dengan
pemikiran logis. Pemaparan latar belakang harus sistematis, logis, serta
disertai data, informasi, dan telaah pustaka dari sumber primer, mutakhir, dan
relevan yang dapat dipertanggungjawabkan secara ilmiah. Masalah penelitian
yang lebih spesifik dirumuskan pada bagian rumusan masalah.

\section{Rumusan Masalah}

Rumusan masalah merupakan pernyataan ringkas mengenai masalah yang akan
diselesaikan dan cara mengatasinya untuk menjawab tujuan tugas akhir. Masalah
yang diteliti dapat dirumuskan karena berbagai sebab, seperti adanya
kesenjangan (\textit{gap}), tantangan, kesangsian, ketidakjelasan, dan
keingintahuan secara akademik yang berkaitan dengan fenomena alam, sosial, dan
ekonomi. Dalam pernyataan ringkas tersebut harus tercakup pendekatan yang
digunakan dalam perumusan masalah.

\section{Tujuan}

Pernyataan tujuan tugas akhir ialah pernyataan singkat dan jelas tentang tujuan
yang akan dicapai sebagai upaya pemecahan masalah maupun memahami gejala
(fenomena) yang dijelaskan dalam latar belakang. Tujuan merupakan pemandu atau
arah untuk merencanakan dan melaksanakan kajian tugas akhir. Tujuan penelitian
ditulis dengan memilih kata kerja yang hasilnya dapat diukur dan dilihat,
seperti: menguraikan, menerangkan, membuktikan, menjajaki, menguji,
membuktikan, atau menerapkan suatu gejala, konsep atau dugaan, atau bahkan
membuat suatu prototipe. Jangan menggunakan kata kerja mengetahui atau
memahami. Masalah dan tujuan penelitian harus terkait dan konsisten.

\section{Manfaat}

Manfaat merupakan dampak positif (kegunaan) dari hasil karya ilmiah tugas akhir
bagi bidang ipteks, pembangunan, dan masyarakat. Manfaat utama dari hasil tugas
akhir adalah menambah khasanah ilmu pengetahuan dalam bentuk pustaka sebagai
sumber acuan/referensi untuk pengembangan ipteks, para pengambil keputusan baik
di industri maupun pemerintah dan lembaga untuk menyusun kebijakan baru, serta
masyarakat umum. Manfaat dinyatakan dengan kata kerja yang lugas dan logis.

\section{Ruang Lingkup (opsional)}

Tugas akhir sering kali dihadapkan pada keterbatasan data, dana, waktu, metode,
bahkan teori. Oleh karena itu, tugas akhir perlu dengan tegas menunjukkan ruang
lingkup dengan mempertimbangkan keterbatasan tersebut.

\section{Hipotesis (opsional)}

Hipotesis dapat ditulis secara eksplisit atau tersirat sesuai bidang ipteks yang
relevan. Hipotesis dapat menjadi bagian dari pendahuluan (untuk bidang sains,
teknik, dan kesehatan), atau merupakan bagian akhir dari tinjauan pustaka
(untuk bidang sosial, ekonomi, dan humaniora).
